\couple{Parentesi Quadrata}{L’ora delle croste 

Si lecca le ferite tutti i giorni alle cinque del pomeriggio, quando si portano a spasso le bestie abbandonate, all’ombra. 
Passa la punta della lingua lungo gli indurimenti e ne stacca gli angoli con i denti, piano, leggermente- non sanguinano. 
Non serve. 

Guaisce a intervalli regolari, scandendo l’ora delle croste.
Presto si arrende e smette di cercare la polpa. Non basta un intervallo pomeridiano per arrivare dalla superficie al midollo, servirebbe del tempo morto che non si trova.
Dopo le ferite si sposta alle congiunzioni tra le dita dei piedi. Si chiede come sia possibile che la pelle, lì, non si sia ancora spaccata, dividendosi. Stira le dita il più possibile fino ai crampi, e poi già si annoia e decide di proseguire oltre.

Si svuota l’ombelico, annusandosi le unghie- odore di camera vecchia.
Lì rimanevano piccoli resti di pelle già usata, che ora diventa pesante; la lascia di fianco a sé, in caso gli venisse fame di chiuso. 

L’incavo delle ascelle lo diverte. Gli piace spingerci un dito dentro fino a sentirne il male in fondo alla gola. Pulsano e odorano, come l’estate sotto il sole dei parcheggi fuori dal centro. 
Cerca di toccarsele con la gola, che non ci arriva, incitando un desiderio felino. 
Tra i peli trova polvere da soffitta, secca. Non la toglie, non nuoce. 

Nell’ora delle croste, quando si portano a spasso le ombre sotto le bestie, si annusa e cerca di scavarsi un minuto, uno soltanto. 
Un minuto fermo, se possibile, che non corra.
Lo scopre nella congiunzione delle sue gambe, dove il sudore si raggruppa e si fa lago, dove i pesci non nuotano, ma stanno a galla. Dove si respira male ma si gioisce meglio.
Nella congiunzione  delle sue gambe, dove i pesci stanno a galla, le bestie abbandonate piangono, mugolano. Lo fanno piano, per paura di essere scoperte e lasciate libere.
 La libertà impedirebbe loro di essere portate a spasso, e a loro preme continuare questa tradizione quotidiana.
Nella congiunzione delle sue gambe la solitudine a volte è condivisa, anche se spesso no. Rimane sola ma comunque piacevole, sudata, unta. 
Alle cinque non gli rimane che quella, questo incrocio che mugola, se aiutato. 

Un minuto fermo per sentire.
}{Scab Time

He licks his wounds every day at five in the afternoon, when the abandoned animals are taken for a walk in the shade. He runs the tip of his tongue along the hardened scab spots and lightly peels away the corners with his teeth, gently. The scabs don't bleed. There is no need to.

He whimpers at regular intervals, marking the time of the scabs.
Soon he gives up and stops searching for the pulp. From the surface to the marrow, an afternoon break isn't enough; it would require downtime, which can't be found.
After the wounds, he moves on to the joints between his toes; he wonders how the skin there hasn't yet split. He stretches the toes apart as hard as it can until they cramp, and then he gets bored and decides to move on.


He empties his bellybutton, and afterwards he sniffs his nails—they smell like an old room.
Small remnants of used skin remained there, now heavy; he leaves them next to him, in case he will get hungry for something old.

The hollows of his armpits amuse him. He likes to push a finger into them until it hurts in the back of his throat. They throb and smell, like the parking lots outside the city center under the summer sun.
He tries to touch them with his throat, but he doesn't reach, inciting a feline desire.
Between the armpit hairs he finds attic dust, dried. He doesn't remove it, because it doesn't harm to leave it there..

In the hour of scabs, when the shadows are taken for a walk beneath the animals, he smells himself and tries to carve out a minute for himself, just one. A minute of stillness, if possible, a minute who doesn’t rush.
He finds it in the junction of his legs, where the sweat becomes a lake, where the fishes don't swim, but float. Where it's hard to breathe but easier to rejoice.
In the junction of his legs where the fishes float, the abandoned animals cry and whimper. They do it quietly, for fear of being discovered and set free. 
They do not want freedom because freedom prevents them from being taken for a walk, and they are anxious to continue the daily tradition.
In the junction of his legs, solitude is sometimes shared, though often not. It remains lonely but still pleasant, sweaty, greasy. At five o'clock, all that's left is this crossroads that whimpers, if helped.

A minute of stillness, to listen.
}

\couple{Raskolnikov}{Cammino per il labirinto composto dal groviglio di vicoli del centro. 
sono le 4:22, ho tardato di 2 minuti il randez-vouz con i miei demoni.
Passo dopo passo, sbuffo dopo sbuffo, un lamento nelle cuffie, un lamento nel vicolo, un lamento nel mio cuore, 3 tossici in un vicolo che spengono la sofferenza con le sostanze.  3 tossici che hanno scelto 3 sipari diversi da calare sulla propria esistenza.  hey, molla questo, ma per fare cosa?
Potevo studiare, ma non me la sento
Potevo scappare, ma non c'era il tempo
Potrei lavorare, ma non ha senso
Preferisco stare qua, stare a fumare,
Solo come un cane, ne ho fumate cento, per le vie del centro, ma non c'è nessuno, solo nero.
piccoli ultimi uomini traviati dalla promessa del fulmine e dalla speranza di essere fulmine, per poi scoprirsi gocce pesanti un istante prima di infrangersi nel fango all' ombra della nube degli ultimi uomini.
non ho nulla, non voglio nulla, non sono nulla. solo fumo e demoni.
Ma allora basta cercare la luce. col favore delle tenebre ribaltiamo questa notte eterna. bruciamo tutto! }{I walk through the maze woven from the tangle of alleys in the old town.
It’s 4:22 a.m., and I’m two minutes late to my rendezvous with my demons.
Step after step, sigh after sigh—
a lament in my headphones,
a lament in the alley,
a lament in my heart.
Three junkies in a side street smother their suffering with chemicals—
three lost souls, each lowering a different curtain on their own existence.
“Hey, drop that,” I whisper, “but for what purpose?”
I could have studied, but I didn’t have it in me.
I could have run, but there was no time.
I could work, I suppose, but what would be the point?
I’d rather just be here, smoking,
alone like a stray dog.
I’ve smoked a hundred of them, wandering through the city center—
yet there’s no one here, only darkness.
We are the small, last men,
led astray by the promise of lightning
and by the hope of becoming lightning,
only to discover ourselves heavy droplets
falling into the mud
beneath the shadowed cloud of the last men.
I have nothing, I want nothing, I am nothing—
nothing but smoke and demons.
So enough with searching for the light.
With the cover of darkness, let’s overturn this eternal night.
Let’s set it all ablaze.}

\couple{Sole Giaguaro}{Un giorno un tizio perse il tatto. Il senso del tatto, intendo. Non la vista, non l’udito. Proprio il tatto.
Si sveglio' una mattina, e nel dormiveglia di una lenta e calda domenica estiva si accorse di non sentirsi addosso le mutande. Per un momento, uno soltanto, non se ne preoccupo’ - beh, era molto caldo, puo’ capitare di dormire nudi - Ma poi, come una valanga. Una sensazione, o meglio, una non-sensazione dopo l’altra lo travolsero: Non sentiva le lenzuola. Non sentiva il cuscino. Non sentiva il materasso sotto di lui! Apri’ gli occhi in preda al panico e si mise a sedere sul letto: tutto era li’ dove doveva essere, il materasso, le lenzuola, pure le sue mutande. Ma non le sentiva. Non le sentiva toccare la sua pelle. Vedeva le lenzuola su di lui, ma non le percepiva. Vedeva le sue mani poggiarsi sul materasso, ma non sentiva nulla.
- Cazzo, sono morto - penso’ - o forse sto ancora sognando.
Fece per darsi un gran pizzicotto, come nei film. Con le dita della mano destra raggiunse il braccio sinistro. Si strinse forte la pelle: non senti’ nulla. Si pizzico' ancora piu’ forte, come se la pelle se la dovesse strappare via. Aveva stretto talmente forte che con le unghie si era graffiato e qualche goccia di sangue aveva iniziato a rigargli la pelle. Stavolta aveva percepito uno strano dolore, lieve, ma in profondita’. Capi’ cosa gli stesse succedendo. Aveva perso il tatto.
Chiamo’ un amico per farsi portare in ospedale. Ci volle una buona mezz’ora per spiegargli quale fosse il problema. Dovette ripetersi almeno un altro paio di volte lungo la strada per l’ospedale, tanto incredibile era la sua condizione. Dopo una lunga giornata passata tra un esame e l’altro alla fine il dottore gli disse che non sapeva che pesci prendere e che, purtroppo, le speranze di trovare un rimedio erano nulle.
Sulla strada del ritorno il suo amico cerco’ di rincuorarlo.
- Non e’ cosi’ terribile dai - diceva l’amico - Alla fine e’ come essere in un videogioco: ti muovi in un mondo virtuale che non puoi toccare realmente.
- Beh, almeno in un videogioco il controller vibra ogni tanto - gli rispose il poveretto - E poi e’ pericoloso: anche il semplice cucinare diventera’ una sfida impossibile.
- Vabe, ma ti preoccupi del cucinare? Per cose del genere devi stare tranquillo, ti aiuteremo noi. Tanto . I cambiamenti piu’ radicali probabilmente saranno altri.
Mentre pronunciava queste parole, una realizzazione li attraverso’ nello stesso momento come un fulmine. Si guardarono l’un l’altro e capirono che stavano pensando alla stessa cosa. Scese il silenzio in macchina. 
Arrivarono a casa. Il suo amico si offri’ di preparargli qualcosa da mangiare. Lui non rispose. Quel pensiero continuava ad ingombrare la sua mente. Un pensiero terribile, una sciagura.
L’amico capiva il suo tormento. Cercava di distrarlo, di chiacchierare con lui. Lui rispondeva a monosillabi. Tutta la sua mente era occupata da quell'agghiacciante pensiero. Sarebbe stato capace di rinunciare a molte cose nella vita: non gli interessava piu’ di tanto di non poter sentire piu’ il caldo o il freddo, non si preoccupava di potersi scottare o di ferirsi facilmente. Ma quello, quello no. Quella realizzazione lo aveva distrutto. Non sapeva se sarebbe riuscito a sostenere un tale peso.
Dopo qualche ora, l’amico lo lascio’. Gli disse di stare tranquillo. Gli disse che magari con una dormita si sarebbe risolto tutto e che domani si sarebbe svegliato come nuovo. Gli disse che molta gente gli voleva bene e che tutti lo avrebbero aiutato. Gli disse di non pensare troppo alle conseguenze, e che sarebbe stato tutto come prima. Mentiva, sapendo di mentire. Era un buon amico. Aveva fatto tutto il possibile.
L’uomo senza tatto, quella notte, si taglio’ le vene dei polsi. Non senti’ quasi dolore. Soltanto un confortevole torpore, quando ormai la maggior parte del suo sangue lo guardava dal pavimento.
Sulla sua lapide, scrissero: “Non pote’ sopportare una vita senza piu' il tocco di una donna”.}{One day a guy lost his touch. The sense of touch, I mean. Not sight, not hearing. Touch itself.
He woke up one morning, and in the half-sleep of a slow and warm summer Sunday he realized he couldn’t feel his underwear on his body. For one moment, just one, he didn’t worry about it - well, it was very hot; it can happen to sleep naked - But then, like an avalanche. One sensation, or rather, one non-sensation after another overwhelmed him: he couldn’t feel the sheets. He couldn’t feel the pillow. He couldn’t feel the mattress beneath him! He opened his eyes in a panic and sat up in bed: everything was where it should be—the mattress, the sheets, even his underwear. But he couldn’t feel them. He couldn’t feel them touching his skin. He saw the sheets on himself, but he didn’t perceive them. He saw his hands resting on the mattress, but he felt nothing.
- Fuck, I’m dead - he thought - or maybe I’m still dreaming.
He tried to give himself a hard pinch, like in the movies. With the fingers of his right hand he grabbed his left arm. He squeezed the skin hard: he felt nothing. He pinched even harder, as if he were about to tear the skin off. He squeezed so hard that he scratched himself with his nails and a few drops of blood began to streak his skin. This time he perceived a strange pain—faint, but deep. He understood what was happening to him. He had lost his sense of touch.
He called a friend to take him to the hospital. It took a good half hour to explain what the problem was. He had to repeat himself at least a couple more times on the way to the hospital, so unbelievable was his condition. After a long day spent going from one test to another, the doctor finally told him he didn’t know what to make of it and that, unfortunately, there was no hope of finding a remedy.
On the way back, his friend tried to comfort him.
- It’s not that terrible - the friend said - In the end it’s like being in a video game: you move around in a virtual world you can’t really touch.
- Well, at least in a video game the controller vibrates now and then - the poor guy replied - And besides, it’s dangerous: even something as simple as cooking will become an impossible challenge.
- Come on, are you worried about cooking now? For things like that you can relax, we’ll help you. Anyway, the most radical changes will probably be others.
As he said these words, a realization struck them both. They looked at each other and understood they were thinking the same thing. A heavy silence fell in the car.
They arrived home. His friend offered to make him something to eat. He didn’t answer. That thought kept cluttering his mind. A terrible thought, a tragedy.
The friend understood his torment. He tried to distract him, to chat with him. He answered in monosyllables. His entire mind was occupied by that chilling thought. He would have been capable of giving up many things in life: he didn’t care much about no longer being able to feel heat or cold, he wasn’t worried about getting burned or injured easily. But that—no. That realization had destroyed him. He didn’t know whether he would be able to bear such a weight.
After a few hours, the friend left him. He told him to relax. He said that maybe after a night’s sleep everything would be resolved and that tomorrow he would wake up like new. He said that many people loved him and that everyone would help him. He told him not to think too much about the consequences, and that everything would be like before. He was lying, knowing he was lying. He was a good friend. He had done everything possible.
The man without touch, that night, slit the veins of his wrists. He felt almost no pain. Only a comforting numbness, when most of his blood was already staring up at him from the floor.
On his gravestone they wrote: “He could not bear a life without the touch of a woman.”}

\couple{Una voce all'orecchio}{Mi sveglio. Penso. Chi sono? Dove sono? Mi riaddormento.
Mi risveglio. Ripenso. Ma è proprio vero che sono? Che fare?
E mi riaddormento.
Apro gli occhi. La luce. Che noia il mondo. Che fare? Non fare. Li richiudo.
Apro gli occhi. Fa freddo. Cosa vuole da me il tutto? Perché? Chi sei?
Corri cazzo, corri! Cosa vuoi? Chi sei?
Sono sempre io.
Corri cazzo, corri!
Dai non anche oggi, perché?
Non senti i fili che ti legano alle cose a cui sei legato?
Non vedi il tempo rotolare verso il basso come una balla di fieno?
No.
Corri cazzo, corri!
Ma dove? Dove?
Non so neanche dove sono.
Eppure corro.
Ciao, come va?
Bene, tu?
Bene, addio.
Corri cazzo, corri!
Il bus. Non so neanche perché abbia le ruote il bus.
Eppure salgo.
Il suo biglietto non è valido per questa corsa.
Ma quale corsa? Ancora sta storia?
Allora ti pago. Stronzo. Io ti pago.
E poi scendo.
Ovviamente la fermata è quella sbagliata.
Come se ce ne fosse una giusta.}{I wake up. I think. Who am I? Where am I? I fall asleep again.
I wake up again. I think again. But is it really true that I am? What should I do?
And I fall asleep again.
I open my eyes. The light. What a bore, the world. What to do? Not do. I close them again.
I open my eyes. It’s cold. What does the whole want from me? Why? Who are you?
Run, fuck, run! What do you want? Who are you?
It’s always me.
Run, fuck, run!
Come on, not today as well, why?
Don’t you feel the strings tying you to the things you are tied to?
Don’t you see time rolling downward like a bale of hay?
No.
Run, fuck, run!
But where? Where?
I don’t even know where I am.
Yet I run.
Hi, how’s it going?
Fine, you?
Fine, goodbye.
Run, fuck, run!
The bus. I don’t even know why the bus has wheels.
Yet I get on.
Your ticket is not valid for this ride.
What ride? This story again?
Fine, I’ll pay you. Asshole. I’ll pay you.
And then I get off.
Obviously it’s the wrong stop.
As if there were a right one.}

\couple{Rupicapra rupicapra}{Suro mi dice di scrivere qualcosa sull’Oblio. 
“Hai sicuramente dei personaggi inerenti!” 
Vorrei trovare una storia che abbia qualcosa da dire. Almeno una chiave di lettura. 
“Penso di avere già in mente quello giusto” 
Mento. Certo, sento di avere un buon bagaglio di personaggi con me. Me li porto dietro come un carretto di burattini. Il problema è che il carretto è disordinato e i miei personaggi non riesco a trovarli. Forse dovrei parlargli. 
“Venite fuori!”
“No…”
“…abbiamo paura. Chi l’ha detto che un personaggio è felice solo quando prende vita. Che ne sai tu di cosa vuol dire vivere per un personaggio, tu che ci tratti come solo come burattini. Piuttosto che fungere da sacchi vuoti, riempiti occasionalmente per le tue storie, vogliamo rimanere qua. Piuttosto che dire cose che non pensiamo, che non conosciamo o comprendiamo preferiamo rimanere muti. Alla luce preferiamo il buio. Dimenticaci, e forse saremo liberi”
“Io vi ho creati! Voi siete quello che siete perché l’ho deciso io! Se c’è una persona che conosce ciò che pensate o sognate quella sono io. Io, che vi conosco meglio di come potrete mai conoscere voi stessi”
“È vero, e proprio per questo non vogliamo uscire nuovamente alla luce. Perché non è una luce che ci rappresenta, è una luce che illumina quello che vuoi tu. Se per conoscere noi stessi il prezzo è il buio, buio sia!” 
Mi sto spazientendo, sento che più questo discorso prosegue e meno ho possibilità di uscirne vincitore. 
“Capisco quello che dite, forse non avete tutti i torti. Però non mi potete abbandonare così, io non vi ho mai voluto fare del male. Senza di voi sono perduto, non saprei più chi sono. Non posso affrontare la vita da solo. Tornerete da me un giorno?”
…nessuna risposta arriva dal loro buio. Sconsolato, chiudo gli occhi. Mi ammutolisco. 
“Perché questo silenzio? Non mi abbandonare adesso ti prego, sto per andare in scena! Senza i tuoi personaggi sono allo sbaraglio! Ci sei? Resta con me!”
È buio, non mi muovo. Nessuna risposta arriva dal mio buio. Volevo trovare una storia con qualcosa da dire. Ma non riesco a dire niente. 

}{Suro tells me to write something about Oblivion.
"You definitely have some relevant characters!"
I'd like to find a story that has something to say. At least a key to understanding it.
"I think I already have the right one in mind."
I'm lying. Sure, I feel like I have a good set of characters with me. I carry them around like a cart of puppets. The problem is that the cart is messy and I can't find my characters. Maybe I should talk to them.
"Come out!"
“No…”
“…we're afraid. Who said a character is only happy when it comes to life? What do you know about what it means to live for a character, you who treat us like mere puppets? Rather than act as empty sacks, occasionally filled for your stories, we want to stay here. Rather than say things we don't think, know, or understand, we prefer to remain silent. We prefer darkness to the light. Forget us, and perhaps we'll be free.”
“I created you! You are what you are because I decided so! If there's anyone who knows what you think or dream, it's me. I, who know you better than you will ever know yourselves.”
“It's true, and that's precisely why we don't want to come out into the light again. Because it's not a light that represents us, it's a light that illuminates what you want. If the price of knowing ourselves is darkness, then so be it!”
I'm getting impatient; I feel like the longer this conversation goes on, the less likely I am to emerge victorious.
“I understand what you're saying; perhaps you're not entirely wrong. But you can't abandon me like this. I never meant to hurt you. Without you, I'm lost; I wouldn't know who I am anymore. I can't face life alone. Will you come back to me someday?”
...no answer comes from their darkness. Dejected, I close my eyes. I fall silent.
“Why this silence? Don't abandon me now, please, I'm about to go on stage! Without your characters, I'm in disarray! Are you there? Stay with me!”
It's dark, I don't move. No answer comes from my darkness. I wanted to find a story with something to say. But I can't say anything.}

